\documentclass[12pt, twoside]{article}
%\documentclass[12pt, twoside]{article}
\usepackage[letterpaper, margin=1in, headsep=0.2in]{geometry}
\setlength{\headheight}{0.6in}
%\usepackage[english]{babel}
\usepackage[utf8]{inputenc}
\usepackage{microtype}
\usepackage{amsmath}
\usepackage{amssymb}
%\usepackage{amsfonts}
\usepackage[nomessages]{fp} %\FPeval{\var-name}{2*sin(pi/6)}
\usepackage{siunitx} %units in math. eg 20\milli\meter
\usepackage{yhmath} % for arcs, overparenth command
\usepackage{tikz} %graphics
\usetikzlibrary{quotes, angles, arrows, arrows.meta}
\usepackage{pgfplots}
\usepgfplotslibrary{fillbetween}
\pgfplotsset{compat=1.16}
\usepackage{graphicx} %consider setting \graphicspath{{images/}}
\usepackage{parskip} %no paragraph indent
\usepackage{enumitem}
\usepackage{multicol}
\usepackage{venndiagram}

\usepackage{fancyhdr}
\pagestyle{fancy}
\fancyhf{}
\renewcommand{\headrulewidth}{0pt} % disable the underline of the header
\raggedbottom
\hfuzz=2mm %suppresses overfull box warnings

\usepackage{hyperref}

\fancyhead[LO]{La Scuola d'Italia / Dr.\ Huson / IB Math 12 \\* 26 October 2026}
\fancyhead[RO]{\fbox{\begin{minipage}{6cm}\footnotesize Nickname:\\[0.5cm]\end{minipage}}}
\fancyfoot[C]{\thepage}

\begin{document}
\subsubsection*{2.4 Pre-Quiz: Sine Rule, Cosine Rule and the Area of a Triangle}

\begin{center}
\fbox{\begin{minipage}{0.9\textwidth}
\centering
$\dfrac{a}{\sin A} = \dfrac{b}{\sin B} = \dfrac{c}{\sin C}
\qquad
c^{2} = a^{2} + b^{2} - 2ab\cos C
\qquad
\text{Area} = \dfrac{1}{2}ab\sin C$
\end{minipage}}
\end{center}

A calculator is allowed. Give answers correct to three significant figures,
unless they are exact. When a question gives an angle in radians, give
angles in radians. Diagrams are not to scale.

\begin{enumerate}[itemsep=0.15cm]

%--- Problem 1: Obtuse angle given (SSA): sine rule for an angle, third angle + sine rule for a side, area [7 marks] ---
%--- Source: Original (freshly authored, AA SL 3.2); modelled on AA SL 2022 Nov P2 Q02 ---
\item[1.]
The following diagram shows triangle $ABC$, with $AC = 21\,\text{cm}$,
$BC = 15\,\text{cm}$ and $A\hat{B}C = 108^\circ$.

\begin{center}
\begin{tikzpicture}[scale=0.3]
  \coordinate (A) at (0,0);
  \coordinate (C) at (21,0);
  \coordinate (B) at (7.904,7.322);
  \draw[thick] (A) -- (B) -- (C) -- cycle;
  \draw pic["$108^\circ$", draw, angle radius=0.5cm, angle eccentricity=1.9] {angle = A--B--C};
  \node[below left] at (A) {$A$};
  \node[above] at (B) {$B$};
  \node[below right] at (C) {$C$};
  \node[below] at (10.5,0) {$21\,\text{cm}$};
  \node[above right] at (14.45,3.66) {$15\,\text{cm}$};
\end{tikzpicture}
\end{center}

\begin{enumerate}[itemsep=0.1cm]
  \item Find $B\hat{A}C$.
  \item Find $AB$.
  \item Find the area of triangle $ABC$.
\end{enumerate}

\begin{tikzpicture}
  \draw (-0.6,0) rectangle (14.6,10.6);
  \foreach \y in {1,2,...,10} \draw[dotted] (0,\y)--(14.1,\y);
\end{tikzpicture}

\newpage

%--- Problem 2: SAS in radians: cosine rule for a side, area, sine rule for an angle in radians [6 marks] ---
%--- Source: Original (freshly authored, AA SL 3.2, 3.4); modelled on AA SL 2022 May P2 TZ2 Q01 ---
\item[2.]
The following diagram shows triangle $PQR$, with $Q\hat{P}R = 1.3$ radians,
$PQ = 9\,\text{cm}$ and $PR = 7\,\text{cm}$.

\begin{center}
\begin{tikzpicture}[scale=0.45]
  \coordinate (P) at (0,0);
  \coordinate (Q) at (9,0);
  \coordinate (R) at (1.869,6.746);
  \draw[thick] (P) -- (Q) -- (R) -- cycle;
  \draw pic["$1.3$", draw, angle radius=0.6cm, angle eccentricity=1.6] {angle = Q--P--R};
  \node[below left] at (P) {$P$};
  \node[below right] at (Q) {$Q$};
  \node[above] at (R) {$R$};
  \node[below] at (4.5,0) {$9\,\text{cm}$};
  \node[left] at (0.93,3.37) {$7\,\text{cm}$};
\end{tikzpicture}
\end{center}

\begin{enumerate}[itemsep=0.1cm]
  \item Find $QR$.
  \item Find the area of triangle $PQR$.
  \item Find $P\hat{Q}R$.
\end{enumerate}

\begin{tikzpicture}
  \draw (-0.6,0) rectangle (14.6,13.0);
  \foreach \y in {1,2,...,12} \draw[dotted] (0,\y)--(14.1,\y);
\end{tikzpicture}

\newpage

%--- Problem 3: Triangular field (SSS): cosine rule for an angle, area, perpendicular distance from area [6 marks] ---
%--- Source: Original (freshly authored, AA SL 3.2, 3.3); modelled on AA SL 2023 May P2 TZ2 Q02 ---
\item[3.]
A farmer owns a triangular field $ABC$, with $AB = 120\,\text{m}$,
$BC = 85\,\text{m}$ and $AC = 140\,\text{m}$. A straight path runs along
the side $AC$.

\begin{center}
\begin{tikzpicture}[scale=0.05]
  \coordinate (A) at (0,0);
  \coordinate (C) at (140,0);
  \coordinate (B) at (95.625,72.497);
  \draw[thick] (A) -- (B) -- (C) -- cycle;
  \node[below left] at (A) {$A$};
  \node[above] at (B) {$B$};
  \node[below right] at (C) {$C$};
  \node[below] at (70,0) {path, $140\,\text{m}$};
  \node[above left] at (47.8,36.2) {$120\,\text{m}$};
  \node[above right] at (117.8,36.2) {$85\,\text{m}$};
\end{tikzpicture}
\end{center}

\begin{enumerate}[itemsep=0.1cm]
  \item Find $A\hat{B}C$.
  \item Find the area of the field.
  \item The farmer walks from $B$ to the path by the shortest possible
  route. Find the length of this route.
\end{enumerate}

\begin{tikzpicture}
  \draw (-0.6,0) rectangle (14.6,12.0);
  \foreach \y in {1,2,...,11} \draw[dotted] (0,\y)--(14.1,\y);
\end{tikzpicture}

\newpage

%--- Problem 4: Ambiguous case in radians: two values of an angle, third angle, area [6 marks] ---
%--- Source: Original (freshly authored, AA SL 3.2); ambiguous case, radian measure ---
\item[4.]
In triangle $ABC$, $AB = 10\,\text{cm}$, $BC = 7\,\text{cm}$ and
$B\hat{A}C = 0.6$ radians.

\begin{enumerate}[itemsep=0.1cm]
  \item Find the two possible values of $A\hat{C}B$.
\end{enumerate}

In the rest of this question, $A\hat{C}B$ is obtuse.

\begin{enumerate}[itemsep=0.1cm]
  \item[(b)] Find $A\hat{B}C$.
  \item[(c)] Find the area of triangle $ABC$.
\end{enumerate}

\begin{tikzpicture}
  \draw (-0.6,0) rectangle (14.6,17.5);
  \foreach \y in {1,2,...,16} \draw[dotted] (0,\y)--(14.1,\y);
\end{tikzpicture}

\newpage

%--- Problem 5: Quadrilateral split by a diagonal: cosine rule for BD, sine rule for an angle, area of quadrilateral [8 marks] ---
%--- Source: Original (freshly authored, AA SL 3.2); IB quadrilateral-with-diagonal pattern ---
\item[5.]
The following diagram shows quadrilateral $ABCD$, with $AB = 6\,\text{cm}$,
$AD = 9\,\text{cm}$, $BC = 8\,\text{cm}$, $B\hat{A}D = 75^\circ$ and
$B\hat{C}D = 62^\circ$.

\begin{center}
\begin{tikzpicture}[scale=0.42]
  \coordinate (A) at (0,0);
  \coordinate (B) at (1.553,5.796);
  \coordinate (C) at (8.363,9.993);
  \coordinate (D) at (9,0);
  \draw[thick] (A) -- (B) -- (C) -- (D) -- cycle;
  \draw[dashed] (B) -- (D);
  \draw pic["$75^\circ$", draw, angle radius=0.5cm, angle eccentricity=1.75] {angle = D--A--B};
  \draw pic["$62^\circ$", draw, angle radius=0.5cm, angle eccentricity=1.8] {angle = B--C--D};
  \node[below left] at (A) {$A$};
  \node[left] at (B) {$B$};
  \node[above] at (C) {$C$};
  \node[below right] at (D) {$D$};
  \node[left] at (0.78,2.9) {$6\,\text{cm}$};
  \node[below] at (4.5,0) {$9\,\text{cm}$};
  \node[above left] at (4.96,7.89) {$8\,\text{cm}$};
\end{tikzpicture}
\end{center}

\begin{enumerate}[itemsep=0.1cm]
  \item Find $BD$.
  \item Find $B\hat{D}C$.
  \item Find the area of quadrilateral $ABCD$.
\end{enumerate}

\begin{tikzpicture}
  \draw (-0.6,0) rectangle (14.6,12.0);
  \foreach \y in {1,2,...,11} \draw[dotted] (0,\y)--(14.1,\y);
\end{tikzpicture}

\newpage

%--- Problem 6: Area equation in x (radians): solve for x, cosine rule for BC, sine rule for an angle in radians [7 marks] ---
%--- Source: Original (freshly authored, AA SL 3.2, 1.x algebra); IB "sides x and x + k, given area" pattern ---
\item[6.]
The following diagram shows triangle $ABC$, with $AB = x\,\text{cm}$,
$AC = (x + 4)\,\text{cm}$ and $B\hat{A}C = 2.1$ radians. The area of
triangle $ABC$ is $30\,\text{cm}^{2}$.

\begin{center}
\begin{tikzpicture}[scale=0.42]
  \coordinate (A) at (0,0);
  \coordinate (B) at (-3.319,5.674);
  \coordinate (C) at (10.574,0);
  \draw[thick] (A) -- (B) -- (C) -- cycle;
  \draw pic["$2.1$", draw, angle radius=0.5cm, angle eccentricity=1.7] {angle = C--A--B};
  \node[below] at (A) {$A$};
  \node[above left] at (B) {$B$};
  \node[below right] at (C) {$C$};
  \node[left] at (-1.66,2.84) {$x$};
  \node[below] at (5.29,0) {$x + 4$};
\end{tikzpicture}
\end{center}

\begin{enumerate}[itemsep=0.1cm]
  \item Find the value of $x$.
  \item Find $BC$.
  \item Find $A\hat{B}C$.
\end{enumerate}

\begin{tikzpicture}
  \draw (-0.6,0) rectangle (14.6,13.7);
  \foreach \y in {1,2,...,13} \draw[dotted] (0,\y)--(14.1,\y);
\end{tikzpicture}

\end{enumerate}
\end{document}
